\documentclass[11pt]{article}
\usepackage[a4paper,margin=27mm]{geometry}
\usepackage[T1]{fontenc}
\usepackage{lmodern}
\usepackage{amsmath,amssymb,amsthm}
\usepackage{microtype}
\usepackage[hidelinks]{hyperref}
\newcommand{\R}{\mathbb R}
\newtheorem{theorem}{Theorem}[section]
\newtheorem{lemma}[theorem]{Lemma}
\newtheorem{corollary}[theorem]{Corollary}
\newtheorem*{informal}{Main result}
\title{Two positive equilibria outside two disks\\for a monostable reaction}
\author{Kaiwen Jin\thanks{Email: \href{mailto:yui.ui.siki@gmail.com}{yui.ui.siki@gmail.com}.}}
\date{}
\begin{document}
\maketitle
\begin{abstract}
We construct a positive monostable reaction and two disjoint disks for which
an exterior reaction--diffusion equation with zero Dirichlet boundary data
has two distinct bounded positive equilibria. This gives a negative answer
to the universal two-obstacle uniqueness question posed by Berestycki and
Graham. One equilibrium lies above a positive profile on the infinite strip
between the disks. A second is obtained by energy minimization on bounded
exhaustions; positive transverse energy bounds its mass in a long narrow
part of the gap. The two mass bounds are incompatible. The reaction and
finite disk geometry are explicit, and both equilibria are classical on
the full exterior domain.
\end{abstract}
\section{Introduction}

Consider bounded positive equilibria of a reaction--diffusion equation
outside absorbing obstacles:
\begin{equation}\label{eq:problem}
 -\Delta u=f(u)\quad\text{in }\Omega,\qquad
 u=0\quad\text{on }\partial\Omega.
\end{equation}
Here the unknown \(u\) is a population density and \(\Omega\) is its habitat.
We assume a positive monostable reaction, in the following precise sense:
\begin{equation}\label{eq:positive-reaction}
\begin{gathered}
 f\in C^{1,\gamma}([0,\infty)),\quad 0<\gamma<1,\qquad f(0)=f(1)=0,\\
 f'(0)>0,\quad f'(1)<0,\qquad
 f>0\text{ on }(0,1),\quad f<0\text{ on }(1,\infty).
\end{gathered}
\end{equation}
There is no assumption that the per-capita growth rate \(f(s)/s\)
decreases. Thus the reaction class is broader than the strong-KPP class
studied in \cite{BGKPP}.

Geometry has a substantial effect on uniqueness in this class.
Berestycki and Graham prove uniqueness outside a compact strongly
star-shaped obstacle, in particular a single ball
\cite[Theorem~1.1]{BG}.
They also explain that two sufficiently separated balls have a unique
positive bounded equilibrium \cite[Section~8.5]{BG}.
Their Open Question~7(ii) asks whether uniqueness holds for the complement
of two balls without restrictions on their radii or separation.
Two nearby balls have no common star center, so the dilation comparison
for a single obstacle does not directly apply. Nonuniqueness for other
domain geometries, including domains with attached pockets and infinite
strips, is already known \cite[Sections~5--6]{BG}. The issue here is the
complement of exactly two balls.

We give a negative answer to the universal uniqueness question.
\begin{informal}
There is one reaction satisfying \eqref{eq:positive-reaction} for every
\(0<\gamma<1\), and two equal disjoint disks in \(\R^2\), for which
\eqref{eq:problem} has at least two distinct bounded positive solutions
in \(C^2(\Omega)\cap C(\overline\Omega)\).
\end{informal}
The reaction, disks and solutions' separating bounds are specified in
Theorem~\ref{thm:main}. One admissible example in dimension two suffices
to refute the assertion quantified over all dimensions, reactions and
disjoint ball pairs. Both solutions are defined on the full exterior
domain and vanish on both complete boundary circles.

The construction uses two features of a pair of large disks with a fixed
small gap. An infinite straight strip fits between them. A positive
one-dimensional profile on this strip, extended by zero, is a subsolution
and produces one equilibrium. At the same time, a long central part of
the actual gap has uniformly bounded width. We choose the reaction so
that the transverse energy is positive for every nonzero profile in this
region. For an equilibrium obtained by energy minimization on expanding
bounded domains, suppressing the population in the middle of the gap
costs energy only in two fixed-length end collars. Its total squared
mass in the gap is therefore bounded independently of the gap's length.
The equilibrium above the strip profile violates this bound when that
length is chosen large.

Section~\ref{sec:reaction} constructs the reaction and proves the
transverse estimate. Section~\ref{sec:geometry} fixes the disks and the
strip profile, and Section~\ref{sec:solutions} constructs and separates
the equilibria. Appendix~\ref{app:limits} gives the regularity, physical
boundary and exhaustion details.
All functions are real valued, integrals are Lebesgue integrals, and
\(H^1_0\) denotes the usual zero-trace Sobolev space.

\section{The reaction and a one-dimensional estimate}\label{sec:reaction}

We first choose a reaction for which a narrow interval has positive energy
for every nonzero profile, although a positive profile can still be a
subsolution. This distinction will produce the two equilibria.

Put
\[
 \theta=\frac14,\qquad h=\frac{13}{50},\qquad \varepsilon=\frac1{100},
 \qquad H(t)=3t^2-2t^3\quad(0\le t\le1),
\]
and define, for \(0\le s\le1\),
\[
 g(s)=
 \begin{cases}
  0,&0\le s\le\theta,\\
  H(100s-25),&\theta<s<h,\\
  1,&h\le s\le3/4,\\
  H(4(1-s)),&3/4<s\le1.
 \end{cases}
\]
Our reaction is
\begin{equation}\label{eq:reaction}
 f(s)=
 \begin{cases}
  g(s)+\varepsilon s(1-s),&0\le s\le1,\\
  -\varepsilon(1-e^{-(s-1)}),&s>1.
 \end{cases}
\end{equation}
The values and first derivatives match at every joining point, since
\(H(0)=0\), \(H(1)=1\), and \(H'(0)=H'(1)=0\).
Moreover, \(f\) and \(f'\) are bounded and \(f'\) is globally Lipschitz.
Consequently \(f\in C^{1,\gamma}([0,\infty))\) for every \(0<\gamma<1\).
It satisfies all the conditions in \eqref{eq:positive-reaction}, with
\(f'(0)=\varepsilon\) and \(f'(1)=-\varepsilon\).

Let \(F(s)=\int_0^s f(t)\,dt\). The estimates needed below are
\begin{equation}\label{eq:primitive}
 0\le f(s)\le2,\qquad F(s)\ge0,\qquad
 F(s)\le(s-\theta)_++\frac{\varepsilon}{2}s^2
 \quad(0\le s\le1).
\end{equation}
Indeed, \(0\le g(s)\le\mathbf1_{\{s>\theta\}}\) and
\(s(1-s)\le s\) on this interval.

\begin{lemma}\label{lem:interval}
Let \(\theta,\mu,\ell>0\) satisfy
\(\mu\ell^2\le16\theta/3\).
Then every nonnegative \(v\in H^1_0(-\ell,\ell)\) satisfies
\begin{equation}\label{eq:interval}
 \frac12\int_{-\ell}^{\ell}|v'|^2
 -\mu\int_{-\ell}^{\ell}(v-\theta)_+\ge0.
\end{equation}
\end{lemma}

\begin{proof}
Use the absolutely continuous representative of \(v\). If
\(E=\{v>\theta\}\) is empty, there is nothing to prove. Otherwise write
\(m=|E|\), so \(0<m<2\ell\). On a component interval \(I\) of \(E\),
the function \(z=v-\theta\) belongs to \(H^1_0(I)\).
For \(I=(a,b)\), set \(p(x)=\mu(x-a)(b-x)/2\), so \(-p''=\mu\) and its endpoint values are zero.
Completing the square gives
\[
 \int_I\left(\frac12|z'|^2-\mu z\right)
 =\frac12\int_I|z'-p'|^2-\frac12\int_I|p'|^2
 \ge-\frac{\mu^2|I|^3}{24}.
\]
Summing over the component intervals, and using
\(\sum_I|I|^3\le m^3\), bounds their total contribution below by
\(-\mu^2m^3/24\).

The two outer intervals, from \(-\ell\) to \(\inf E\) and from
\(\sup E\) to \(\ell\), have positive lengths \(L_-\) and \(L_+\).
Their endpoint values are respectively \(0,\theta\) and \(\theta,0\).
By Cauchy--Schwarz their gradient energy is at least
\[
 \frac{\theta^2}{2}\left(\frac1{L_-}+\frac1{L_+}\right)
 \ge\frac{2\theta^2}{L_-+L_+}
 \ge\frac{2\theta^2}{2\ell-m}.
\]
All remaining gradient terms are nonnegative. Thus the left-hand side of
\eqref{eq:interval} is at least
\[
 \frac{2\theta^2}{2\ell-m}-\frac{\mu^2m^3}{24}.
\]
This is nonnegative: the maximum of \(m^3(L-m)\) on \([0,L]\) is
\(27L^4/256\), attained at \(m=3L/4\), and the hypothesis with
\(L=2\ell\) is exactly the required inequality.
The argument also covers countably many components, since the integrals
of both \(|v'|^2\) and \((v-\theta)_+\) are finite.
\end{proof}

The following consequence will control the population in the gap. Its
width bound is fixed now so that the same estimate applies to every
horizontal section of the region constructed in Section~\ref{sec:geometry}.

\begin{corollary}\label{cor:coercive}
For the reaction \eqref{eq:reaction}, let \(\ell_1=9/8\).
If \(0<\ell\le\ell_1\) and \(v\in H^1_0(-\ell,\ell)\) takes values in
\([0,1]\), then
\begin{equation}\label{eq:coercive}
 \int_{-\ell}^{\ell}\left(\frac12|v'|^2-F(v)\right)
 \ge\frac1{40}\int_{-\ell}^{\ell}v^2.
\end{equation}
\end{corollary}

\begin{proof}
Apply Lemma~\ref{lem:interval} with \(\mu=26/25\) and
\(\theta=1/4\). Its hypothesis holds because
\(\mu\ell_1^2=1053/800<4/3\).
Together with \eqref{eq:primitive}, this gives
\[
 \int\left(\frac12|v'|^2-F(v)\right)
 \ge\frac1{52}\int|v'|^2-\frac1{200}\int v^2.
\]
The interval Poincar\'e constant is \(\pi^2/(4\ell^2)\).
For example, for compactly supported smooth \(v\), write
\(v=\varphi z\) with \(\varphi(x)=\cos(\pi x/(2\ell))\), and integrate
by parts to obtain
\[
 \int|v'|^2-\frac{\pi^2}{4\ell^2}\int v^2
 =\int\varphi^2|z'|^2\ge0.
\]
Density gives the same inequality on \(H^1_0\). Since
\(\pi^2/(4\ell^2)\ge16/9\), we conclude that
\[
 \int\left(\frac12|v'|^2-F(v)\right)
 \ge\left(\frac4{117}-\frac1{200}\right)\int v^2
 =\frac{683}{23400}\int v^2
 \ge\frac1{40}\int v^2.
\]
\end{proof}

\section{Two disks and a strip subsolution}\label{sec:geometry}

Take
\begin{equation}\label{eq:constants}
 \ell_0=\frac{11}{10},\qquad T=20000,\qquad
 R=40(T+2)^2=16\,003\,200\,160.
\end{equation}
Let
\begin{equation}\label{eq:domain}
 c_\pm=\bigl(\pm(R+\ell_0),0\bigr),\qquad
 \Omega=\R^2\setminus
 \bigl(\overline{B_R(c_+)}\cup\overline{B_R(c_-)}\bigr).
\end{equation}
The disks are disjoint and the distance between them is \(2\ell_0=11/5\).
The large radius makes a long part of their gap narrow enough for
Corollary~\ref{cor:coercive}. More precisely, set
\[
 k(y)=\ell_0+R-\sqrt{R^2-y^2}\quad(|y|<R),\qquad
 D=\{(x,y):|y|<T,\ |x|<k(y)\}.
\]
The lateral sides of \(D\) are portions of the two physical circles.
For \(|y|\le T+2\),
\begin{equation}\label{eq:width}
 \ell_0\le k(y)
 =\ell_0+\frac{y^2}{R+\sqrt{R^2-y^2}}
 \le\ell_0+\frac{y^2}{R}\le\ell_1.
\end{equation}

\begin{theorem}\label{thm:main}
For the reaction \(f\) in \eqref{eq:reaction} and the two-disk exterior
domain \(\Omega\) in \eqref{eq:domain}, problem \eqref{eq:problem} has
two distinct solutions \(u,U\in C^2(\Omega)\cap C(\overline\Omega)\).
Both satisfy \(0<u,U<1\) throughout \(\Omega\), and
\begin{equation}\label{eq:separation}
 \int_Du^2\,dx\,dy\le4500,\qquad
 \int_DU^2\,dx\,dy>5000.
\end{equation}
\end{theorem}

All parameters in this theorem are finite and fixed. In particular,
the theorem is an example in dimension two with strictly separated
obstacles; it does not use a limiting domain of parallel walls.
The particular numerical values are chosen for convenient estimates,
without optimizing the radii or gap.

To construct \(U\), we use the entire strip
\(\{|x|<\ell_0\}\), which lies inside \(\Omega\). Recall that
\(h=13/50\) and \(f(s)\ge1\) for \(h\le s\le3/4\). Put
\[
 q=\frac{11+\sqrt{17}}{20},\qquad
 A=h+\frac{q^2}{2}=\frac{173+11\sqrt{17}}{400},
\]
and define
\begin{equation}\label{eq:w}
 w(x)=
 \begin{cases}
  A-x^2/2,&|x|\le q,\\
  q(\ell_0-|x|),&q<|x|<\ell_0,\\
  0,&|x|\ge\ell_0.
 \end{cases}
\end{equation}
Since \(q(\ell_0-q)=h\), both \(w\) and \(w'\) agree at
\(x=\pm q\). On \(|x|<q\), we have \(h\le w\le A<3/4\),
so \(f(w)\ge1=-w''\). On the linear pieces \(-w''=0\le f(w)\).
At each of \(x=\pm\ell_0\), the derivative jumps upward by \(q\);
therefore \(-w''\) has a negative Dirac mass there. It follows that
\begin{equation}\label{eq:subsolution}
 W(x,y):=w(x),\qquad -\Delta W\le f(W)
\end{equation}
in the sense of distributions on the whole plane, and hence on
\(\Omega\). On both physical circles \(W=0\), because their points
have \(|x|\ge\ell_0\).

We will also use the strict estimate
\begin{equation}\label{eq:strip-mass}
 \int_{-\ell_0}^{\ell_0}w(x)^2\,dx>\frac18.
\end{equation}
Indeed, \(q>3/4\) and \(A>217/400\), so
\(w(x)\ge A-1/32>1/2\) when \(|x|\le1/4\).

\section{Construction of the two equilibria}\label{sec:solutions}

We work first on bounded domains
\[
 \Omega_n=B_n(0)\cap\Omega,\qquad
 n\in\mathbb N,\quad n>2R+\ell_0+2,
\]
with boundary values zero on the two disks and one on the outer circle.
These data have the \(H^1\) extension
\(b(z)=\min\{1,\operatorname{dist}(z,\Omega^c)\}\).
Only the outer radius \(n\) will tend to infinity. The reaction,
the obstacles and the region \(D\) stay fixed.

We use the following bounded-domain Sobolev facts: bounded \(H^1\)
sequences on a bounded Lipschitz domain admit weak \(H^1\) and strong
\(L^2\) subsequences; traces are continuous; and zero-trace functions
are approximable in \(H^1\) by smooth compactly supported functions.
These facts, the truncation rules and Poincar\'e's inequality are
recalled in \cite[Section~1.1]{FRRO}.
The interior regularity and the preservation of the physical boundary
values in our exhaustion are detailed in Appendix~\ref{app:limits}.

\subsection{An equilibrium above the strip profile}

The functions \(W\) and \(1\) are an ordered subsolution and
supersolution on every \(\Omega_n\), compatible with the prescribed
boundary values. Here is the bounded-domain construction. Choose
\(K\ge\max\{1,\operatorname{Lip}_{[0,1]}f\}\), set \(v_0=W\), and solve
\[
 (-\Delta+K)v_{j+1}=f(v_j)+Kv_j
\]
with the fixed boundary values. The linear Dirichlet problem is solved
by minimizing
\[
 \frac12\int_{\Omega_n}(|\nabla v|^2+Kv^2)
 -\int_{\Omega_n}(f(v_j)+Kv_j)v
\]
in \(b+H^1_0(\Omega_n)\). This quadratic is coercive and strictly convex. Its comparison principle follows by testing a
weak supersolution against its negative part.
Since \(s\mapsto f(s)+Ks\) is nondecreasing on \([0,1]\), comparison
gives
\[
 W\le v_j\le v_{j+1}\le1.
\]
The right-hand sides are uniformly bounded. Testing the linear equations
against \(v_j-b\) and using Poincar\'e's inequality bounds \(v_j\) in
\(H^1(\Omega_n)\). Weak compactness, strong \(L^2\) convergence and
bounded convergence then give a weak solution \(U_n\) with
\(W\le U_n\le1\). It is classical in the interior.

Lemma~\ref{lem:limits} yields a subsequential exterior limit \(U\)
with the complete zero boundary values and \(U\ge W\).
As \(W\not\equiv0\), Lemma~\ref{lem:strict} gives \(0<U<1\).

\subsection{A minimizing exhaustion}

For this construction only, extend \(f\) oddly to the negative half-line
and let \(F\) denote its primitive on \(\R\).
The extension is \(C^1\), bounded, and has bounded derivative.
Its primitive is even, globally Lipschitz, and maximized at
\(\pm1\).
On \(\Omega_n\) minimize
\[
 J_n(v)=\int_{\Omega_n}
 \left(\frac12|\nabla v|^2-F(v)\right)\,dx\,dy
 \quad\text{over }b+H^1_0(\Omega_n).
\]
The direct method applies. Since \(F\) is bounded above,
a minimizing sequence has bounded gradients; Poincar\'e's inequality
then bounds it in \(H^1\).
For a strongly \(L^2\)-convergent subsequence, the potential integrals
converge by
\[
 \left|\int_{\Omega_n}F(v_j)-F(v)\right|
 \le\operatorname{Lip}(F)|\Omega_n|^{1/2}
       \|v_j-v\|_{L^2(\Omega_n)}.
\]
The gradient energy is weakly lower semicontinuous and the trace is
preserved, so a global minimizer exists.

Replacing it by \(\min\{1,|v|\}\) preserves the boundary data and cannot
increase the energy: the gradient decreases and the value of \(F\)
increases. We may therefore choose a minimizer \(u_n\) with
\(0\le u_n\le1\). It is still a minimizer in the \emph{unrestricted}
affine class. Thus
\[
 -\Delta u_n=f(u_n)
\]
weakly, and its second variation gives
\begin{equation}\label{eq:stable}
 \int_{\Omega_n}\bigl(|\nabla\phi|^2-f'(u_n)\phi^2\bigr)\,dx\,dy
 \ge0\qquad(\phi\in C_c^\infty(\Omega_n)).
\end{equation}
The differentiations are justified by the boundedness of \(f\) and
\(f'\). Interior regularity makes \(u_n\) classical.
We next obtain a bound in the long gap that is uniform in \(n\).

\subsection{An energy bound in the gap}

Let \(\eta(y)\) be zero for \(|y|\le T-1\), one for \(|y|\ge T\),
and linear in each of the two intervening unit intervals.
Then \(|\eta'|\le1\). Inside \(D\), replace \(u_n\) by \(\eta u_n\),
leaving it unchanged elsewhere. This is a legal variation:
\((\eta-1)u_n\) has zero trace on the lateral physical boundary,
and on the horizontal cuts \(y=\pm T\) because \(\eta=1\).
Its zero extension belongs to \(H^1_0(\Omega_n)\).
Minimality and \eqref{eq:primitive} imply
\begin{equation}\label{eq:remove}
 J_D(u_n)\le J_D(\eta u_n)
 \le\int_C(|\nabla u_n|^2+1)\,dx\,dy,
 \qquad C=\{(x,y)\in D:T-1<|y|<T\}.
\end{equation}
Here \(J_D\) denotes the same energy restricted to \(D\).
For the last inequality we used \(F(\eta u_n)\ge0\) and
\(\tfrac12|\eta\nabla u_n+u_n\eta'e_y|^2
 \le|\nabla u_n|^2+1\).

To bound the gradient in these collars, take a Lipschitz function
\(\chi(y)\) equal to one for \(T-1\le|y|\le T\), zero outside
\(T-2<|y|<T+1\), and linear on the remaining two ramps at each end.
It satisfies \(|\chi'|\le1\).
Use \(\chi^2u_n\) as a test on the central channel
\(\{|y|<T+2,\ |x|<k(y)\}\), and extend it by zero.
It has zero trace on the physical sides and near the horizontal ends,
so it belongs to \(H^1_0(\Omega_n)\).
The weak equation and Young's inequality give
\begin{align}
 \int\chi^2|\nabla u_n|^2
 &\le2\int\chi^2 f(u_n)u_n
       +4\int u_n^2|\chi'|^2 \notag\\
 &\le8\cdot\frac{27}{2}=108. \label{eq:caccioppoli}
\end{align}
In fact, the support has total \(y\)-length \(6\) and width at most
\(2\ell_1=9/4\), hence area at most \(27/2\); also \(f(u_n)u_n\le2\).
The area of \(C\) is at most \(9/2\). From \eqref{eq:remove},
\begin{equation}\label{eq:neck-energy}
 J_D(u_n)\le108+\frac92=\frac{225}{2}.
\end{equation}

For almost every \(y\in(-T,T)\), the horizontal section of \(u_n\)
belongs to \(H^1_0(-k(y),k(y))\): its \(x\)-derivative is square
integrable by Fubini, and its endpoint values on the physical circles
are zero by Lemma~\ref{lem:limits}.
Applying Corollary~\ref{cor:coercive} on each such section and retaining
the nonnegative \(y\)-gradient contribution, we obtain
\[
 J_D(u_n)\ge\frac1{40}\int_Du_n^2\,dx\,dy.
\]
Consequently
\begin{equation}\label{eq:low-mass}
 \int_Du_n^2\,dx\,dy\le4500.
\end{equation}

\subsection{Proof of Theorem~\ref{thm:main}}

By Lemma~\ref{lem:limits}, a subsequence of \(u_n\) converges locally
in \(C^2(\Omega)\) to a solution \(u\) with the complete zero boundary
values. The stability inequalities \eqref{eq:stable} and
Lemma~\ref{lem:nonzero} exclude the zero limit.
Lemma~\ref{lem:strict} therefore gives \(0<u<1\).
Since \(D\) has finite area and \(0\le u_n\le1\), dominated convergence
in \eqref{eq:low-mass} yields
\[
 \int_Du^2\,dx\,dy\le4500.
\]
On the other hand, \(U\ge W\), and every horizontal section of \(D\)
contains \((-\ell_0,\ell_0)\). By \eqref{eq:strip-mass},
\[
 \int_DU^2\,dx\,dy
 \ge2T\int_{-\ell_0}^{\ell_0}w(x)^2\,dx
 >\frac{T}{4}=5000.
\]
Thus \(u\ne U\), and both are solutions of the same exterior problem
in the required class. This proves the theorem.

\appendix
\section{Regularity, physical boundary values and exterior limits}
\label{app:limits}

We give the analytic details used in the bounded-domain constructions.
They also make explicit why the limiting functions solve the problem on
the whole exterior, including both complete physical boundaries.

\begin{lemma}\label{lem:limits}
For the fixed reaction \eqref{eq:reaction} and fixed domain
\eqref{eq:domain}, let \(v_n\in H^1(\Omega_n)\) solve
\(-\Delta v_n=f(v_n)\), take values in \([0,1]\), and have zero trace
on the two inner circles.
There is a subsequence converging in \(C^2_{\mathrm{loc}}(\Omega)\)
to a solution \(v\in C^2(\Omega)\cap C(\overline\Omega)\), with
\(0\le v\le1\) and zero values on both full circles.
\end{lemma}

\begin{proof}
We first justify the interior estimates for initially weak solutions.
On a fixed ball \(B_{2r}\Subset\Omega\), write
\[
 P(z)=\frac1{2\pi}\int_{B_{2r}}
       \log\frac1{|z-\zeta|}\,f(v_n(\zeta))\,d\zeta.
\]
The logarithmic kernel and its first derivative are locally integrable
in two dimensions. Thus \(P\) and \(\nabla P\) are uniformly bounded
on fixed bounded sets, and
\(-\Delta P=\mathbf1_{B_{2r}}f(v_n)\) distributionally.
The coefficient \(1/(2\pi)\) follows from the flux identity
\(-\int_{\partial B_\rho}\partial_\nu((2\pi)^{-1}
\log(1/|z|))=1\).
After a compact smooth cutoff, Morrey's inequality at \(p=4\)
gives a uniform \(C^{0,1/2}\) bound for \(P\).
The remainder \(v_n-P\) is weakly harmonic and bounded, hence smooth
with uniform interior derivative bounds.
Here we use \cite[Theorem~1.5 and Corollary~1.12]{FRRO}.

It follows that \(v_n\), and then \(f(v_n)\), have uniform local
\(C^{0,1/2}\) bounds. Mollify on a smaller ball by a nonnegative
smooth kernel of integral one. The smooth mollifications satisfy the
Poisson equation with uniformly \(C^{0,1/2}\) right-hand sides.
The interior Schauder estimate
\[
 \|a\|_{C^{2,\alpha}(B_{1/2})}
 \le C_\alpha\bigl(\|a\|_{L^\infty(B_1)}
                    +\|\Delta a\|_{C^{0,\alpha}(B_1)}\bigr),
 \qquad 0<\alpha<1,
\]
for smooth \(a\), is \cite[Theorem~2.12]{FRRO}.
Translate and rescale it to the fixed balls, with \(\alpha=1/2\).
Letting the mollification scale tend to zero proves uniform local
\(C^{2,1/2}\) bounds for \(v_n\).
Compactness on fixed balls and a countable diagonal selection now give
a locally \(C^2\)-convergent subsequence, whose limit solves the
same equation.

For the boundary, choose \(\delta=1/2\). The annulus
\(R<|z-c_\pm|<R+\delta\) lies in \(\Omega\), since the gap between
the two disks is \(11/5>\delta\).
It lies in every \(\Omega_n\) considered here.
On this annulus set
\[
 B(r)=-\frac{r^2-R^2}{2}+C_R\log(r/R),\qquad
 C_R=\frac{1+((R+\delta)^2-R^2)/2}{\log((R+\delta)/R)}.
\]
Then \(-\Delta B=2\), \(B(R)=0\) and \(B(R+\delta)=1\).
The weak maximum principle gives \(B\ge0\).
Since
\(-\Delta(B-v_n)=2-f(v_n)\ge0\) and the boundary traces of
\(B-v_n\) are nonnegative, the same principle gives \(0\le v_n\le B\).
This is precisely the weak comparison in
\cite[Proposition~1.13]{FRRO}.
The bound is uniform in \(n\), and \(B(r)\to0\) as \(r\downarrow R\).
It therefore gives the continuous zero value of the limit at every
point of both circles.
\end{proof}

\begin{lemma}\label{lem:nonzero}
If, in addition, the functions \(v_n\) in Lemma~\ref{lem:limits}
satisfy
\[
 \int_{\Omega_n}(|\nabla\phi|^2-f'(v_n)\phi^2)\ge0
 \qquad(\phi\in C_c^\infty(\Omega_n)),
\]
then their exterior limit is not identically zero.
\end{lemma}

\begin{proof}
For each fixed compactly supported test \(\phi\), pass to the limit in
this inequality. If \(v\equiv0\), then \(f'(0)=\varepsilon>0\) would give
\[
 \int_\Omega|\nabla\phi|^2\ge
 \varepsilon\int_\Omega\phi^2
 \qquad(\phi\in C_c^\infty(\Omega)).
\]
Fix a nonzero smooth function \(\psi\) supported in the unit ball.
For each \(S>0\), the ball centered at \((0,4R+4S)\) with radius
\(S\) lies in the fixed exterior domain. The translated dilation of
\(\psi\) into this ball has Rayleigh quotient
\[
 S^{-2}\frac{\int|\nabla\psi|^2}{\int\psi^2}.
\]
Taking \(S\) large contradicts the preceding inequality, since
\(\varepsilon>0\).
\end{proof}

\begin{lemma}\label{lem:strict}
Every nonzero solution \(v\) furnished by Lemma~\ref{lem:limits}
satisfies \(0<v<1\) in \(\Omega\).
\end{lemma}

\begin{proof}
Since \(-\Delta v=f(v)\ge0\), the strong minimum principle gives
\(v>0\). In this setting the principle follows directly from spherical
means: for a nonnegative superharmonic \(C^2\) function, the spherical
average about a point with value zero is at most zero; hence the function
vanishes on a ball, and then on its connected domain.

For the upper bound put \(z=1-v\ge0\), and take
\(L\ge\max\{1,\operatorname{Lip}_{[0,1]}f\}\).
Then \(\Delta z=f(v)\le Lz\).
On
\(\Omega\times(-\pi/(2\sqrt L),\pi/(2\sqrt L))\), the function
\(z(x,y)\cos(\sqrt L\,t)\) is nonnegative and superharmonic.
An interior zero would make it identically zero by the same
minimum principle, so \(v\equiv1\).
This contradicts its continuous zero values on the physical boundary.
\end{proof}

\section*{AI assistance}
OpenAI Codex was used in the mathematical development, source checking
and drafting of this manuscript.

\begin{thebibliography}{9}
\bibitem{BG}
H.~Berestycki and C.~Graham,
\emph{A Stable-Compact Method for Qualitative Properties of Semilinear
Elliptic Equations},
Archive for Rational Mechanics and Analysis \textbf{250} (2026), article~21.
\url{https://doi.org/10.1007/s00205-026-02168-6}.

\bibitem{BGKPP}
H.~Berestycki and C.~Graham,
\emph{The steady states of strong-KPP reactions in general domains},
arXiv:2212.06611 (2022).
\url{https://arxiv.org/abs/2212.06611}.

\bibitem{FRRO}
X.~Fern\'andez-Real and X.~Ros-Oton,
\emph{Regularity Theory for Elliptic PDE},
author-hosted version, preface dated Z\"urich 2020.
\url{https://user.math.uzh.ch/ros-oton/Llibre-ellipticPDE.pdf}.
The theorem and page locators used here refer to this version.
Published by EMS Press, 2022, DOI
\url{https://doi.org/10.4171/ZLAM/28}.
\end{thebibliography}

\end{document}
